\documentclass[english,twoside, openright, hidelinks, 11 pt, class=report,crop=false]{standalone}
\input{../../preamb}
\input{../bm}

\begin{document}	
\vedlegg{Eksakte verdier for utvalgte vinkler} \label{sincoseksakt}
\begin{center}
	\renewcommand{\arraystretch}{1.5}	
	\begin{tabular}{l|c|c|c|c|c|c|c|c|c}
		& 0&$30^\circ$ & $45^\circ$ &$60^\circ$ & $90^\circ$ & $ 120^\circ $ & $ 135^\circ$ & $ 130^\circ $ & $ 180^\circ $    \\
		\hline
		$\sin x$ & 0 &$\frac{1}{2}$ & $\frac{\sqrt{2}}{2}$ & $\frac{\sqrt{3}}{2}$ & 1 & $ \frac{\sqrt{3}}{2} $ & $ \frac{\sqrt{2}}{2} $ & $ \frac{1}{2} $ & $ 0 $\\
		$\cos x$ & 1 & $\frac{\sqrt{3}}{2}$ & $\frac{\sqrt{2}}{2}$ & $\frac{1}{2}$ & 0 & $ -\frac{1}{2} $ & $ -\frac{\sqrt{2}}{2} $ & $ -\frac{\sqrt{3}}{2} $ & $ -1 $
	\end{tabular}
\end{center}
\subsubsection{Tips}
For å finne eksakte verdier av sinus til en vinkel, kan vi sette opp følgende tabell:\vs \renewcommand{\arraystretch}{1.5}	
\begin{center}
	\begin{tabular}{l|c|c|c|c|c|}
		& $0^\circ$&$30^\circ$ & $45^\circ$ &$60^\circ$ & $90^\circ$    \\
		\hline
		$\sin x$ & 0 &$\frac{\sqrt{1}}{2}$ & $\frac{\sqrt{2}}{2}$ & $\frac{\sqrt{3}}{2}$ & $ \frac{\sqrt{4}}{2} $ \\
	\end{tabular}
\end{center}
For cosinus setter vi opp mønsteret andre veien:
\begin{center}
	\begin{tabular}{l|c|c|c|c|c}
		$\cos x$ & $\frac{\sqrt{4}}{2}  $ & $\frac{\sqrt{3}}{2}$ & $\frac{\sqrt{2}}{2}$ & $\frac{\sqrt{1}}{2}$ & 0 \\
	\end{tabular}
\end{center}
Erstatter vi $ {\frac{\sqrt{1}}{2}} $ med $ \frac{1}{2} $ og $ \frac{\sqrt{4}}{2} $ med 1, får vi dette:

\begin{center}
	\renewcommand{\arraystretch}{1.5}	
	\begin{tabular}{l|c|c|c|c|c}
		& $0^\circ$&$30^\circ$ & $45^\circ$ &$60^\circ$ & $90^\circ$    \\
		\hline
		$\sin x$ & 0 &$\frac{1}{2}$ & $\frac{\sqrt{2}}{2}$ & $\frac{\sqrt{3}}{2}$ & 1 \\
		$\cos x$ & 1 & $\frac{\sqrt{3}}{2}$ & $\frac{\sqrt{2}}{2}$ & $\frac{1}{2}$ & 0 \\
	\end{tabular}
\end{center}
Av tabellen over kan vi enkelt finne $\tan x= \frac{\sin x}{\cos x} $:
\begin{center}
	\renewcommand{\arraystretch}{1.5}	
	\begin{tabular}{l|c|c|c|c|c}
		& $0^\circ$&$30^\circ$ & $45^\circ$ &$60^\circ$ & $90^\circ$    \\
		\hline
		$\sin x$ & 0 &$\frac{1}{2}$ & $\frac{\sqrt{2}}{2}$ & $\frac{\sqrt{3}}{2}$ & 1 \\
		$\cos x$ & 1 & $\frac{\sqrt{3}}{2}$ & $\frac{\sqrt{2}}{2}$ & $\frac{1}{2}$ & 0 \\
		$\tan x$ & 0 &$\frac{1}{\sqrt{3}}$ & $1$ & $\sqrt{3}$ & $ \infty  $
	\end{tabular}
\end{center}


\end{document}